\begin{bt}%[9D6H18-2]
% nguon: Đề tham khảo Tuyển sinh 10 TP.HCM 2025-2026 – Quận 1 – Đề 1 – Bài 1
Cho parabol $(P): y = \dfrac{x^2}{2}$.
\begin{enumerate}[a)]
	\item Vẽ đồ thị $(P)$ của hàm số trên.
	\item Tìm trên $(P)$ những điểm $M$ sao cho tung độ gấp đôi hoành độ.
\end{enumerate}
\loigiai{
\begin{enumerate}[a)]
	\item Bảng giá trị tương ứng giữa $x$ và $y$ của hàm số $y = \dfrac{x^2}{2}$:
	\begin{center}
	\begin{tabular}{|c|c|c|c|c|c|}
	\hline
	$x$ & $-4$ & $-2$ & $0$ & $2$ & $4$ \\
	\hline
	$y = \dfrac{x^2}{2}$ & $8$ & $2$ & $0$ & $2$ & $8$ \\
	\hline
	\end{tabular}
	\end{center}
	Đồ thị $(P)$ là một đường parabol đi qua gốc toạ độ $O(0; 0)$, nhận trục tung $Oy$ làm trục đối xứng và nằm phía trên trục hoành.
	\item Điểm $M(x; y)$ thuộc $(P)$ có tung độ gấp đôi hoành độ thoả mãn:
	$y = 2x \Leftrightarrow \dfrac{x^2}{2} = 2x \Leftrightarrow x^2 - 4x = 0 \Leftrightarrow x(x - 4) = 0 \Leftrightarrow \left[\begin{aligned} x &= 0 \Rightarrow y = 0 \\ x &= 4 \Rightarrow y = 8. \end{aligned}\right.$\\
	Vậy có hai điểm thoả mãn yêu cầu đề bài là $M_1(0; 0)$ và $M_2(4; 8)$.
\end{enumerate}
}
\end{bt}

\begin{bt}%[9D6H20-2]
% nguon: Đề tham khảo Tuyển sinh 10 TP.HCM 2025-2026 – Quận 1 – Đề 1 – Bài 2
Cho phương trình $2x^2 - 9x + 3 = 0$.
\begin{enumerate}[a)]
	\item Chứng tỏ phương trình trên có hai nghiệm phân biệt.
	\item Gọi $x_1, x_2$ là hai nghiệm của phương trình. Tính giá trị của biểu thức:
	$M = x_1(2025 - x_1) + x_2(2026 - x_2) - x_2$.
\end{enumerate}
\loigiai{
\begin{enumerate}[a)]
	\item Xét phương trình $2x^2 - 9x + 3 = 0$, ta có:
	$\Delta = b^2 - 4ac = (-9)^2 - 4\cdot 2\cdot 3 = 81 - 24 = 57 > 0$.\\
	Vì $\Delta > 0$ nên phương trình luôn có hai nghiệm phân biệt $x_1, x_2$.
	\item Theo định lí Viète, ta có:
	$\begin{cases} x_1 + x_2 = -\dfrac{b}{a} = \dfrac{9}{2} \\ x_1 x_2 = \dfrac{c}{a} = \dfrac{3}{2}. \end{cases}$\\
	Biến đổi biểu thức $M$:
	$M = 2025x_1 - x_1^2 + 2026x_2 - x_2^2 - x_2 = 2025(x_1 + x_2) - (x_1^2 + x_2^2) = 2025(x_1 + x_2) - \left[(x_1 + x_2)^2 - 2x_1 x_2\right]$.\\
	Thay $x_1 + x_2 = \dfrac{9}{2}$ và $x_1 x_2 = \dfrac{3}{2}$ vào biểu thức $M$, ta được:
	$M = 2025\cdot\dfrac{9}{2} - \left[\left(\dfrac{9}{2}\right)^2 - 2\cdot\dfrac{3}{2}\right] = \dfrac{18225}{2} - \left(\dfrac{81}{4} - 3\right) = \dfrac{18225}{2} - \dfrac{69}{4} = \dfrac{36381}{4} = 9095{,}25$.\\
	Vậy $M = \dfrac{36381}{4} = 9095{,}25$.
\end{enumerate}
}
\end{bt}

\begin{bt}%[9D8H26-2]
% nguon: Đề tham khảo Tuyển sinh 10 TP.HCM 2025-2026 – Quận 1 – Đề 1 – Bài 3
Một chiếc hộp chứa 3 viên bi gồm: một viên bi màu xanh (X), một viên bi màu vàng (V) và một viên bi màu đỏ (Đ). Các viên bi có cùng kích thước và khối lượng. Bạn An lần lượt lấy ra ngẫu nhiên từng viên bi từ trong hộp cho đến khi hộp hết viên bi.
\begin{enumerate}[a)]
	\item Xác định không gian mẫu của phép thử.
	\item Tính xác suất của mỗi biến cố sau:
	\begin{itemize}
		\item Biến cố $A$: ``Viên bi màu đỏ được lấy ra đầu tiên''.
		\item Biến cố $B$: ``Viên bi màu xanh được lấy ra trước viên bi màu vàng''.
		\item Biến cố $C$: ``Viên bi lấy ra lần cuối cùng không có màu xanh''.
	\end{itemize}
\end{enumerate}
\loigiai{
\begin{enumerate}[a)]
	\item Không gian mẫu $\Omega$ là tập hợp tất cả các thứ tự lấy ra 3 viên bi từ hộp:
	$\Omega = \{(\text{X}, \text{V}, \text{Đ}); (\text{X}, \text{Đ}, \text{V}); (\text{V}, \text{X}, \text{Đ}); (\text{V}, \text{Đ}, \text{X}); (\text{Đ}, \text{X}, \text{V}); (\text{Đ}, \text{V}, \text{X})\}$.\\
	Số phần tử của không gian mẫu là $n(\Omega) = 3! = 6$.
	\item Xác suất của từng biến cố:
	\begin{itemize}
		\item Các kết quả thuận lợi cho biến cố $A$ là: $(\text{Đ}, \text{X}, \text{V})$ và $(\text{Đ}, \text{V}, \text{X}) \Rightarrow n(A) = 2$.\\
		Do đó $P(A) = \dfrac{n(A)}{n(\Omega)} = \dfrac{2}{6} = \dfrac{1}{3}$.
		\item Các kết quả thuận lợi cho biến cố $B$ (bi màu xanh xuất hiện trước bi màu vàng) là: $(\text{X}, \text{V}, \text{Đ})$, $(\text{X}, \text{Đ}, \text{V})$ và $(\text{Đ}, \text{X}, \text{V}) \Rightarrow n(B) = 3$.\\
		Do đó $P(B) = \dfrac{n(B)}{n(\Omega)} = \dfrac{3}{6} = \dfrac{1}{2}$.
		\item Biến cố $C$ xảy ra khi viên bi cuối cùng là màu vàng hoặc màu đỏ. Các kết quả thuận lợi cho biến cố $C$ là: $(\text{X}, \text{V}, \text{Đ})$, $(\text{X}, \text{Đ}, \text{V})$, $(\text{V}, \text{X}, \text{Đ})$ và $(\text{Đ}, \text{X}, \text{V}) \Rightarrow n(C) = 4$.\\
		Do đó $P(C) = \dfrac{n(C)}{n(\Omega)} = \dfrac{4}{6} = \dfrac{2}{3}$.
	\end{itemize}
\end{enumerate}
}
\end{bt}

\begin{bt}%[9D6V21-3]
% nguon: Đề tham khảo Tuyển sinh 10 TP.HCM 2025-2026 – Quận 1 – Đề 1 – Bài 4
Một chiếc bàn tròn có bán kính là $2\text{ m}$. Người ta muốn thiết kế một chiếc khăn trải bàn hình tròn có bán kính rộng hơn bán kính của mặt bàn là $x\text{ (m)}$ ($x > 0$).
\begin{enumerate}[a)]
	\item Viết biểu thức $S$ biểu diễn diện tích của tấm khăn trải bàn theo $x$.
	\item Biết rằng phần diện tích khăn rủ xuống ngoài mép bàn là $12{,}0576\text{ m}^2$. Tính giá trị của $x$ (lấy $\pi \approx 3{,}14$).
\end{enumerate}
\loigiai{
\begin{enumerate}[a)]
	\item Bán kính của tấm khăn trải bàn hình tròn là: $R = 2 + x\text{ (m)}$.\\
	Biểu thức biểu diễn diện tích của tấm khăn trải bàn là:
	$S = \pi R^2 = \pi (2 + x)^2 = \pi (x^2 + 4x + 4)\text{ (m}^2\text{)}$.
	\item Diện tích mặt bàn tròn là: $S_{\text{bàn}} = \pi \cdot 2^2 = 4\pi\text{ (m}^2\text{)}$.\\
	Phần diện tích khăn rủ xuống ngoài mép bàn là:
	$S_{\text{dư}} = S - S_{\text{bàn}} = \pi (x^2 + 4x + 4) - 4\pi = \pi (x^2 + 4x)\text{ (m}^2\text{)}$.\\
	Theo giả thiết, $S_{\text{dư}} = 12{,}0576\text{ m}^2$ và $\pi \approx 3{,}14$, ta có:
	$3{,}14(x^2 + 4x) = 12{,}0576 \Leftrightarrow x^2 + 4x = 3{,}84 \Leftrightarrow x^2 + 4x - 3{,}84 = 0$.\\
	Ta có $\Delta' = 2^2 - 1\cdot(-3{,}84) = 4 + 3{,}84 = 7{,}84 = (2{,}8)^2 > 0$.\\
	Vì $x > 0$ nên phương trình có nghiệm dương duy nhất:
	$x = -2 + 2{,}8 = 0{,}8\text{ (m)}$.\\
	Vậy giá trị của $x$ là $0{,}8\text{ mét}$.
\end{enumerate}
}
\end{bt}

\begin{bt}%[9H0V31-3]
% nguon: Đề tham khảo Tuyển sinh 10 TP.HCM 2025-2026 – Quận 1 – Đề 1 – Bài 5
Bạn Linh rất thích trà sữa trân châu của một tiệm nổi tiếng.
\begin{enumerate}[a)]
	\item Một ly trà sữa có dạng hình trụ size S có bán kính đáy $R = 40\text{ mm} = 4\text{ cm}$, chiều cao ly là $h = 160\text{ mm} = 16\text{ cm}$. Tiệm rót $40\%$ thể tích ly là trà sữa, thêm $35\%$ thể tích là đá viên. Linh đoán rằng ly trà sữa cho vào khoảng $10$ viên trân châu dạng hình cầu đường kính mỗi viên $10\text{ mm} = 1\text{ cm}$ sẽ vừa khít và không bị tràn trà sữa ra ngoài. Linh đoán vậy có đúng không? Em hãy giải thích.
	\item Nếu Linh chọn lấy lượng đá là $50\%$ thể tích ly (giữ nguyên $40\%$ thể tích trà sữa), hỏi ly trà sữa đó chứa được tối đa bao nhiêu viên trân châu để không bị tràn ra ngoài? (Lấy $\pi \approx 3{,}14$).
\end{enumerate}
\loigiai{
\begin{enumerate}[a)]
	\item Thể tích của chiếc ly hình trụ là:
	$V_{\text{ly}} = \pi R^2 h = 3{,}14 \cdot 4^2 \cdot 16 = 803{,}84\text{ (cm}^3\text{)}$.\\
	Tổng thể tích trà sữa và đá viên đã chiếm:
	$40\% + 35\% = 75\%$ thể tích ly.\\
	Thể tích phần còn trống trong ly là:
	$V_{\text{trống}} = (100\% - 75\%) \cdot V_{\text{ly}} = 25\% \cdot 803{,}84 = 200{,}96\text{ (cm}^3\text{)}$.\\
	Mỗi viên trân châu hình cầu có bán kính $r = \dfrac{1}{2} = 0{,}5\text{ (cm)}$. Thể tích một viên trân châu là:
	$V_1 = \dfrac{4}{3}\pi r^3 = \dfrac{4}{3} \cdot 3{,}14 \cdot (0{,}5)^3 \approx 0{,}523\text{ (cm}^3\text{)}$.\\
	Thể tích của $10$ viên trân châu là:
	$V_{10} = 10 \cdot V_1 \approx 10 \cdot 0{,}523 = 5{,}23\text{ (cm}^3\text{)}$.\\
	Vì $V_{10} = 5{,}23\text{ cm}^3$ nhỏ hơn rất nhiều so với phần thể tích còn trống $200{,}96\text{ cm}^3$ (còn trống tới gần $195{,}7\text{ cm}^3$), nên ly trà sữa không thể ``vừa khít''. Do đó Linh đoán không đúng.
	\item Khi lấy $50\%$ thể tích là đá và $40\%$ thể tích là trà sữa thì chất lỏng và đá chiếm $90\%$ thể tích ly.\\
	Thể tích phần còn trống lúc này là:
	$V'_{\text{trống}} = 10\% \cdot V_{\text{ly}} = 10\% \cdot 803{,}84 = 80{,}384\text{ (cm}^3\text{)}$.\\
	Số viên trân châu tối đa có thể cho vào ly để không bị tràn là:
	$N = \dfrac{V'_{\text{trống}}}{V_1} = \dfrac{80{,}384}{\frac{4}{3}\cdot 3{,}14\cdot 0{,}125} \approx \dfrac{80{,}384}{0{,}5233} \approx 153{,}6$.\\
	Vì số viên trân châu là số nguyên nên ly có thể chứa tối đa $153$ viên trân châu để không tràn ra ngoài.
\end{enumerate}
}
\end{bt}

\begin{bt}%[9D1V3-4]
% nguon: Đề tham khảo Tuyển sinh 10 TP.HCM 2025-2026 – Quận 1 – Đề 1 – Bài 6
Anh Bình đang có kế hoạch tiết kiệm để mua một chiếc xe ô tô cũ có giá $250\text{ triệu đồng}$ trong vòng một năm ($12\text{ tháng}$). Hiện tại, anh Bình đã có $50\text{ triệu đồng}$ trong tài khoản tiết kiệm.
\begin{enumerate}[a)]
	\item Hãy xác định số tiền tối thiểu anh Bình cần tiết kiệm mỗi tháng để đạt được mục tiêu mua xe sau $12\text{ tháng}$.
	\item Nếu vào tháng thứ $7$, anh Bình phải dùng $15\text{ triệu đồng}$ từ khoản tiết kiệm để gửi về hỗ trợ gia đình, thì trung bình mỗi tháng trong $12\text{ tháng}$ anh Bình cần tiết kiệm ít nhất bao nhiêu tiền để vẫn đạt được mục tiêu mua xe?
\end{enumerate}
\loigiai{
\begin{enumerate}[a)]
	\item Số tiền anh Bình còn thiếu để đủ $250\text{ triệu đồng}$ mua xe là:
	$250 - 50 = 200\text{ (triệu đồng)}$.\\
	Số tiền tối thiểu anh Bình cần tiết kiệm mỗi tháng trong $12\text{ tháng}$ là:
	$\dfrac{200}{12} = \dfrac{50}{3}\text{ (triệu đồng)} \approx 16{,}67\text{ triệu đồng} = 16\,666\,667\text{ đồng}$.
	\item Khi rút $15\text{ triệu đồng}$ gửi về cho gia đình, tổng số tiền anh Bình cần tiết kiệm trong cả năm là:
	$200 + 15 = 215\text{ (triệu đồng)}$.\\
	Số tiền ít nhất anh Bình cần tiết kiệm trung bình mỗi tháng là:
	$\dfrac{215}{12}\text{ (triệu đồng)} \approx 17{,}92\text{ triệu đồng} = 17\,916\,667\text{ đồng}$.
\end{enumerate}
}
\end{bt}

\begin{bt}%[9H9C29-3]
% nguon: Đề tham khảo Tuyển sinh 10 TP.HCM 2025-2026 – Quận 1 – Đề 1 – Bài 7
Từ điểm $M$ nằm ngoài đường tròn $(O; R)$, kẻ hai tiếp tuyến $MA, MB$ ($A, B$ là hai tiếp điểm) và cát tuyến $MCD$ ($C$ nằm giữa $M$ và $D$, $AC > BC$). Gọi $E$ là trung điểm của dây $CD$.
\begin{enumerate}[a)]
	\item Chứng minh $OE \perp CD$ tại $E$ và tứ giác $AOEB$ nội tiếp một đường tròn.
	\item Gọi $F$ là giao điểm của $MD$ và $AB$. Chứng minh $EM$ là tia phân giác của góc $\widehat{AEB}$ và $EF^2 = EA \cdot EB - FA \cdot FB$.
	\item Trong trường hợp $OM = 2R$ và $R = 10\text{ cm}$, tính số đo góc $\widehat{AOB}$ và tính diện tích phần chung của hình tròn $(O)$ và hình tròn đường kính $OM$ (làm tròn kết quả đến hàng phần mười của $\text{cm}^2$).
\end{enumerate}
\loigiai{
\begin{enumerate}[a)]
	\item Vì $E$ là trung điểm của dây $CD$ không đi qua tâm $O$ nên $OE \perp CD$ tại $E$, hay $\widehat{OEM} = 90^\circ$.\\
	Vì $MA, MB$ là các tiếp tuyến của $(O)$ nên $MA \perp OA \Rightarrow \widehat{OAM} = 90^\circ$ và $MB \perp OB \Rightarrow \widehat{OBM} = 90^\circ$.\\
	Năm điểm $O, A, E, B, M$ cùng thuộc đường tròn đường kính $OM$ vì các tam giác vuông $\Delta OAM, \Delta OEM, \Delta OBM$ đều có cạnh huyền là $OM$.\\
	Do đó bốn điểm $A, O, E, B$ cùng thuộc một đường tròn, hay tứ giác $AOEB$ nội tiếp.
	\item Trong đường tròn đường kính $OM$, ta có $MA = MB$ (tính chất hai tiếp tuyến cắt nhau), nên hai cung $MA$ và $MB$ bằng nhau.\\
	Suy ra hai góc nội tiếp $\widehat{AEM}$ và $\widehat{BEM}$ chắn hai cung bằng nhau, do đó $\widehat{AEM} = \widehat{BEM}$.\\
	Vậy $EM$ (hay $EF$) là tia phân giác của góc $\widehat{AEB}$.\\
	Áp dụng tính chất đường phân giác trong tam giác $EAB$, ta có công thức độ dài đường phân giác:
	$EF^2 = EA \cdot EB - FA \cdot FB$.
	\item Trong tam giác vuông $OAM$ tại $A$, ta có:
	$\cos\widehat{AOM} = \dfrac{OA}{OM} = \dfrac{R}{2R} = \dfrac{1}{2} \Rightarrow \widehat{AOM} = 60^\circ$.\\
	Vì $OM$ là tia phân giác của góc $\widehat{AOB}$ nên:
	$\widehat{AOB} = 2\widehat{AOM} = 2\cdot 60^\circ = 120^\circ$.\\
	Gọi $(I)$ là đường tròn đường kính $OM$. Bán kính của $(I)$ là $R_I = \dfrac{OM}{2} = R = 10\text{ cm}$.\\
	Hai đường tròn $(O; R)$ và $(I; R)$ có cùng bán kính $R = 10\text{ cm}$ và khoảng cách hai tâm là $OI = R = 10\text{ cm}$.\\
	Dây cung chung $AB$ cắt đoạn nối tâm $OI$ tại trung điểm $H$ của $OI$ với $OH = \dfrac{R}{2}$.\\
	Trong $(O)$, tam giác $OAB$ cân tại $O$ có $\widehat{AOB} = 120^\circ$. Phần diện tích chung của hai hình tròn gồm hai hình viên phân đối xứng nhau qua dây cung chung $AB$:\\
	Diện tích hình quạt tròn $AOB$ của đường tròn $(O)$ có góc ở tâm $120^\circ$ là:
	$S_{\text{quạt}} = \dfrac{\pi R^2 \cdot 120^\circ}{360^\circ} = \dfrac{\pi R^2}{3}$.\\
	Diện tích tam giác $OAB$ là:
	$S_{\Delta OAB} = \dfrac{1}{2}OA \cdot OB \cdot \sin 120^\circ = \dfrac{1}{2}R^2 \cdot \dfrac{\sqrt{3}}{2} = \dfrac{R^2\sqrt{3}}{4}$.\\
	Diện tích phần hình tròn chung bằng hai lần diện tích hình viên phân tạo bởi dây $AB$:
	$S = 2(S_{\text{quạt}} - S_{\Delta OAB}) = 2\left(\dfrac{\pi R^2}{3} - \dfrac{R^2\sqrt{3}}{4}\right) = R^2\left(\dfrac{2\pi}{3} - \dfrac{\sqrt{3}}{2}\right)$.\\
	Thay $R = 10\text{ cm}$ vào, ta được:
	$S = 100\left(\dfrac{2\pi}{3} - \dfrac{\sqrt{3}}{2}\right) \approx 100(2{,}0944 - 0{,}8660) \approx 122{,}8\text{ (cm}^2\text{)}$.\\
	Vậy diện tích phần chung xấp xỉ bằng $122{,}8\text{ cm}^2$.
\end{enumerate}
}
\end{bt}

\begin{bt}%[9D6H18-2]
% nguon: Đề tham khảo Tuyển sinh 10 TP.HCM 2025-2026 – Quận 1 – Đề 2 – Bài 1
Cho parabol $(P): y = -\dfrac{1}{2}x^2$.
\begin{enumerate}[a)]
	\item Vẽ đồ thị $(P)$ trên mặt phẳng toạ độ $Oxy$.
	\item Tìm các điểm $M$ thuộc đồ thị $(P)$ sao cho $M$ có hoành độ và tung độ là hai số đối nhau.
\end{enumerate}
\loigiai{
\begin{enumerate}[a)]
	\item Bảng giá trị tương ứng giữa $x$ và $y$ của hàm số $y = -\dfrac{1}{2}x^2$:
	\begin{center}
	\begin{tabular}{|c|c|c|c|c|c|}
	\hline
	$x$ & $-4$ & $-2$ & $0$ & $2$ & $4$ \\
	\hline
	$y = -\dfrac{1}{2}x^2$ & $-8$ & $-2$ & $0$ & $-2$ & $-8$ \\
	\hline
	\end{tabular}
	\end{center}
	Đồ thị $(P)$ là một đường parabol có đỉnh là gốc toạ độ $O(0; 0)$, nhận trục tung $Oy$ làm trục đối xứng và nằm phía dưới trục hoành.
	\item Điểm $M(x; y)$ thuộc $(P)$ có hoành độ và tung độ là hai số đối nhau thoả mãn:
	$y = -x \Leftrightarrow -\dfrac{1}{2}x^2 = -x \Leftrightarrow x^2 - 2x = 0 \Leftrightarrow x(x - 2) = 0 \Leftrightarrow \left[\begin{aligned} x &= 0 \Rightarrow y = 0 \\ x &= 2 \Rightarrow y = -2. \end{aligned}\right.$\\
	Vậy có hai điểm thoả mãn là $M_1(0; 0)$ và $M_2(2; -2)$.
\end{enumerate}
}
\end{bt}

\begin{bt}%[9D6H20-2]
% nguon: Đề tham khảo Tuyển sinh 10 TP.HCM 2025-2026 – Quận 1 – Đề 2 – Bài 2
Cho phương trình $3x^2 - 2x - 2 = 0$.
\begin{enumerate}[a)]
	\item Chứng tỏ phương trình trên có hai nghiệm phân biệt.
	\item Không giải phương trình, hãy tính giá trị của biểu thức $T = \dfrac{x_1}{x_2 - 1} + \dfrac{x_2}{x_1 - 1}$.
\end{enumerate}
\loigiai{
\begin{enumerate}[a)]
	\item Xét phương trình $3x^2 - 2x - 2 = 0$, ta có:
	$a\cdot c = 3\cdot(-2) = -6 < 0$.\\
	Vì tích $ac < 0$ nên phương trình luôn có hai nghiệm phân biệt trái dấu $x_1, x_2$.
	\item Theo định lí Viète, ta có:
	$\begin{cases} x_1 + x_2 = -\dfrac{b}{a} = \dfrac{2}{3} \\ x_1 x_2 = \dfrac{c}{a} = -\dfrac{2}{3}. \end{cases}$\\
	Quy đồng mẫu số biểu thức $T$:
	$T = \dfrac{x_1(x_1 - 1) + x_2(x_2 - 1)}{(x_2 - 1)(x_1 - 1)} = \dfrac{x_1^2 + x_2^2 - (x_1 + x_2)}{x_1 x_2 - (x_1 + x_2) + 1} = \dfrac{(x_1 + x_2)^2 - 2x_1 x_2 - (x_1 + x_2)}{x_1 x_2 - (x_1 + x_2) + 1}$.\\
	Thay $x_1 + x_2 = \dfrac{2}{3}$ và $x_1 x_2 = -\dfrac{2}{3}$ vào biểu thức $T$, ta được:
	$T = \dfrac{\left(\dfrac{2}{3}\right)^2 - 2\cdot\left(-\dfrac{2}{3}\right) - \dfrac{2}{3}}{-\dfrac{2}{3} - \dfrac{2}{3} + 1} = \dfrac{\dfrac{4}{9} + \dfrac{4}{3} - \dfrac{2}{3}}{-\dfrac{4}{3} + 1} = \dfrac{\dfrac{4}{9} + \dfrac{6}{9}}{-\dfrac{1}{3}} = \dfrac{\dfrac{10}{9}}{-\dfrac{1}{3}} = -\dfrac{10}{3}$.\\
	Vậy $T = -\dfrac{10}{3}$.
\end{enumerate}
}
\end{bt}

\begin{bt}%[9D8H26-2]
% nguon: Đề tham khảo Tuyển sinh 10 TP.HCM 2025-2026 – Quận 1 – Đề 2 – Bài 3
Một chiếc hộp có chứa $5$ tấm thẻ cùng loại, được đánh số lần lượt là $3; 5; 6; 7; 9$. Lấy ra ngẫu nhiên đồng thời $2$ tấm thẻ từ hộp.
\begin{enumerate}[a)]
	\item Xác định không gian mẫu và số kết quả có thể xảy ra của phép thử.
	\item Tính xác suất của biến cố $A$: ``Tích các số ghi trên $2$ tấm thẻ chia hết cho $3$''.
\end{enumerate}
\loigiai{
\begin{enumerate}[a)]
	\item Không gian mẫu $\Omega$ gồm tất cả các tập hợp con gồm $2$ thẻ chọn từ $5$ thẻ:
	$\Omega = \{(3; 5); (3; 6); (3; 7); (3; 9); (5; 6); (5; 7); (5; 9); (6; 7); (6; 9); (7; 9)\}$.\\
	Số kết quả có thể xảy ra của phép thử là $n(\Omega) = 10$.
	\item Tích của hai số chia hết cho $3$ khi và chỉ khi ít nhất một trong hai thẻ mang số chia hết cho $3$.\\
	Trong $5$ thẻ đã cho, các số chia hết cho $3$ là: $3; 6; 9$ (gồm $3$ thẻ); các số không chia hết cho $3$ là: $5; 7$ (gồm $2$ thẻ).\\
	Cặp duy nhất có tích không chia hết cho $3$ là $(5; 7)$ (vì $5\cdot 7 = 35$).\\
	Do đó số kết quả thuận lợi cho biến cố $A$ là: $n(A) = 10 - 1 = 9$.\\
	Xác suất của biến cố $A$ là:
	$P(A) = \dfrac{n(A)}{n(\Omega)} = \dfrac{9}{10} = 0{,}9$.
\end{enumerate}
}
\end{bt}

\begin{bt}%[9D6V21-3]
% nguon: Đề tham khảo Tuyển sinh 10 TP.HCM 2025-2026 – Quận 1 – Đề 2 – Bài 4
Một mảnh đất hình chữ nhật có chiều dài là $3x + 10\text{ (m)}$ và chiều rộng là $2x + 5\text{ (m)}$. Ở giữa mảnh đất người ta đào một cái ao hình tam giác đều có cạnh là $x + 1\text{ (m)}$, phần đất còn lại dùng để trồng rau.
\begin{enumerate}[a)]
	\item Viết biểu thức $A$ biểu thị diện tích phần đất trồng rau.
	\item Nếu diện tích của cái ao là $12\sqrt{3}\text{ m}^2$ thì diện tích phần trồng rau là bao nhiêu mét vuông (làm tròn kết quả đến hàng đơn vị)?
\end{enumerate}
\loigiai{
\begin{enumerate}[a)]
	\item Diện tích mảnh đất hình chữ nhật là:
	$S_{\text{đất}} = (3x + 10)(2x + 5) = 6x^2 + 15x + 20x + 50 = 6x^2 + 35x + 50\text{ (m}^2\text{)}$.\\
	Diện tích cái ao hình tam giác đều cạnh $x + 1$ là:
	$S_{\text{ao}} = \dfrac{\sqrt{3}}{4}(x + 1)^2\text{ (m}^2\text{)}$.\\
	Biểu thức $A$ biểu thị diện tích phần đất trồng rau là:
	$A = S_{\text{đất}} - S_{\text{ao}} = 6x^2 + 35x + 50 - \dfrac{\sqrt{3}}{4}(x + 1)^2\text{ (m}^2\text{)}$.
	\item Theo đề bài, diện tích cái ao là $12\sqrt{3}\text{ m}^2$, nên ta có:
	$\dfrac{\sqrt{3}}{4}(x + 1)^2 = 12\sqrt{3} \Leftrightarrow (x + 1)^2 = 48 \Leftrightarrow x + 1 = \sqrt{48} = 4\sqrt{3} \Leftrightarrow x = 4\sqrt{3} - 1 \approx 5{,}928\text{ (m)}$.\\
	Diện tích mảnh đất hình chữ nhật lúc này là:
	$S_{\text{đất}} = 6(4\sqrt{3} - 1)^2 + 35(4\sqrt{3} - 1) + 50 = 6(49 - 8\sqrt{3}) + 140\sqrt{3} - 35 + 50 = 294 - 48\sqrt{3} + 140\sqrt{3} + 15 = 309 + 92\sqrt{3}\text{ (m}^2\text{)}$.\\
	Diện tích phần trồng rau là:
	$A = S_{\text{đất}} - S_{\text{ao}} = (309 + 92\sqrt{3}) - 12\sqrt{3} = 309 + 80\sqrt{3} \approx 309 + 138{,}564 = 447{,}564\text{ (m}^2\text{)}$.\\
	Làm tròn đến hàng đơn vị, diện tích phần trồng rau xấp xỉ bằng $448\text{ m}^2$.
\end{enumerate}
}
\end{bt}

\begin{bt}%[9H0V31-3]
% nguon: Đề tham khảo Tuyển sinh 10 TP.HCM 2025-2026 – Quận 1 – Đề 2 – Bài 5
Một khối kim loại đặc bằng sắt có dạng hình hộp chữ nhật với kích thước ba cạnh là $12\text{ cm}, 10\text{ cm}, 7\text{ cm}$ bị khoét một rãnh suốt chiều dài $12\text{ cm}$ bởi một nửa hình trụ có đường kính đáy $4\text{ cm}$.
\begin{enumerate}[a)]
	\item Tính thể tích của khối kim loại còn lại (làm tròn kết quả đến hàng đơn vị của $\text{cm}^3$).
	\item Nấu nóng chảy toàn bộ khối kim loại còn lại này để đúc các viên bi sắt hình cầu có bán kính $1\text{ cm}$. Hỏi có thể đúc được nhiều nhất bao nhiêu viên bi sắt như thế? (Giả sử hao hụt khi đúc không đáng kể, lấy $\pi \approx 3{,}14$).
\end{enumerate}
\loigiai{
\begin{enumerate}[a)]
	\item Thể tích khối hộp chữ nhật ban đầu là:
	$V_{\text{hộp}} = 12 \cdot 10 \cdot 7 = 840\text{ (cm}^3\text{)}$.\\
	Nửa hình trụ bị khoét có bán kính đáy $R = \dfrac{4}{2} = 2\text{ (cm)}$ và chiều cao $h = 12\text{ cm}$.\\
	Thể tích phần kim loại bị khoét đi là:
	$V_{\text{khoét}} = \dfrac{1}{2}\pi R^2 h = \dfrac{1}{2} \cdot 3{,}14 \cdot 2^2 \cdot 12 = 24\cdot 3{,}14 = 75{,}36\text{ (cm}^3\text{)}$.\\
	Thể tích của khối kim loại còn lại là:
	$V_{\text{còn}} = V_{\text{hộp}} - V_{\text{khoét}} = 840 - 75{,}36 = 764{,}64\text{ (cm}^3\text{)}$.\\
	Làm tròn đến hàng đơn vị, thể tích khối kim loại còn lại xấp xỉ bằng $765\text{ cm}^3$.
	\item Thể tích của một viên bi sắt hình cầu bán kính $r = 1\text{ cm}$ là:
	$V_{\text{bi}} = \dfrac{4}{3}\pi r^3 = \dfrac{4}{3} \cdot 3{,}14 \cdot 1^3 \approx 4{,}187\text{ (cm}^3\text{)}$.\\
	Số viên bi sắt tối đa có thể đúc được là:
	$n = \dfrac{V_{\text{còn}}}{V_{\text{bi}}} = \dfrac{764{,}64}{4{,}187} \approx 182{,}6$.\\
	Vì số viên bi phải là số nguyên dương nên có thể đúc được nhiều nhất $182$ viên bi sắt.
\end{enumerate}
}
\end{bt}

\begin{bt}%[9D1V3-4]
% nguon: Đề tham khảo Tuyển sinh 10 TP.HCM 2025-2026 – Quận 1 – Đề 2 – Bài 6
Chị Ba muốn mở một cửa hàng bán trà sữa với giá niêm yết là $24\,000\text{ đồng/ly}$. Nhân dịp khai trương, chị Ba cân nhắc hai hình thức khuyến mãi:
\begin{itemize}
	\item Phương án 1: Giảm giá $30\%$ cho mỗi ly trà sữa.
	\item Phương án 2: Mua $2$ ly được tặng thêm $1$ ly cùng loại.
\end{itemize}
Theo em, nếu hiệu ứng thu hút khách hàng của hai phương án là như nhau thì chị Ba nên chọn phương án nào để có lợi hơn cho cửa hàng? Vì sao?
\loigiai{
Xét số tiền thu được của cửa hàng đối với từng phương án trên mỗi đơn vị sản phẩm bán ra:
\begin{itemize}
	\item Theo Phương án 1 (giảm giá $30\%$ mỗi ly):\\
	Giá bán thực tế mỗi ly trà sữa là:
	$24\,000 \cdot (100\% - 30\%) = 24\,000 \cdot 0{,}7 = 16\,800\text{ (đồng/ly)}$.\\
	Do đó, mỗi ly bán ra cửa hàng thu về $16\,800\text{ đồng}$.
	\item Theo Phương án 2 (mua $2$ tặng $1$):\\
	Khi khách hàng nhận $3$ ly trà sữa thì chỉ phải trả tiền $2$ ly. Số tiền khách trả là:
	$2 \cdot 24\,000 = 48\,000\text{ (đồng)}$.\\
	Giá bán trung bình mỗi ly trà sữa mà cửa hàng giao cho khách là:
	$\dfrac{48\,000}{3} = 16\,000\text{ (đồng/ly)}$.\\
	Do đó, mỗi ly xuất ra cửa hàng chỉ thu về $16\,000\text{ đồng}$.
\end{itemize}
Vì $16\,800\text{ đồng} > 16\,000\text{ đồng}$ (chênh lệch $800\text{ đồng}$ trên mỗi ly) nên Phương án 1 mang lại doanh thu trên từng ly cao hơn so với Phương án 2.\\
Vậy chị Ba nên chọn Phương án 1 để có lợi hơn cho cửa hàng.
}
\end{bt}

\begin{bt}%[9H9C29-3]
% nguon: Đề tham khảo Tuyển sinh 10 TP.HCM 2025-2026 – Quận 1 – Đề 2 – Bài 7
Cho tam giác $ABC$ có ba góc nhọn ($AB < AC$), đường tròn tâm $(O)$ đường kính $BC$ cắt $AB, AC$ lần lượt tại $D$ và $E$. Đoạn thẳng $BE$ cắt $CD$ tại $H$; tia $AH$ cắt $BC$ tại $F$.
\begin{enumerate}[a)]
	\item Chứng minh $AF \perp BC$ tại $F$ và $\widehat{HEF} = \widehat{HCF}$.
	\item Gọi $K$ là giao điểm của $ED$ và $BC$. Chứng minh $EB$ là tia phân giác của góc $\widehat{DEF}$ và $FO \cdot FK = FB \cdot FC$.
	\item Tiếp tuyến tại $B$ của $(O)$ cắt đường thẳng $KE$ tại $I$. Gọi $J$ là trung điểm của đoạn thẳng $AH$. Chứng minh $OI \perp BJ$.
\end{enumerate}
\loigiai{
\begin{enumerate}[a)]
	\item Điểm $D, E$ thuộc đường tròn $(O)$ đường kính $BC$ nên $\widehat{BDC} = \widehat{BEC} = 90^\circ$ (góc nội tiếp chắn nửa đường tròn).\\
	Suy ra $CD \perp AB$ và $BE \perp AC$. Do đó $CD$ và $BE$ là hai đường cao của tam giác $ABC$.\\
	Vì $BE$ cắt $CD$ tại $H$ nên $H$ là trực tâm tam giác $ABC$, suy ra $AH \perp BC$ tại $F$, tức là $AF \perp BC$.\\
	Xét tứ giác $CEHF$, ta có $\widehat{CEH} = 90^\circ$ và $\widehat{CFH} = 90^\circ$ (do $AF \perp BC$).\\
	Hai đỉnh kề $E$ và $F$ cùng nhìn đoạn thẳng $CH$ dưới góc $90^\circ$ nên $CEHF$ là tứ giác nội tiếp đường tròn đường kính $CH$.\\
	Suy ra $\widehat{HEF} = \widehat{HCF}$ (hai góc nội tiếp cùng chắn cung $HF$).
	\item Ta có tứ giác $BDHF$ nội tiếp (do $\widehat{BDH} + \widehat{BFH} = 90^\circ + 90^\circ = 180^\circ$) nên $\widehat{HDF} = \widehat{HBF}$.\\
	Mặt khác, trong $(O)$, ta có $\widehat{HBF} = \widehat{EBC} = \widehat{EDC}$ (cùng chắn cung $EC$), hay $\widehat{HBF} = \widehat{HDE}$.\\
	Suy ra $\widehat{HDE} = \widehat{HDF}$, do đó $DH$ (hay $DC$) là phân giác của $\widehat{EDF}$.\\
	Chứng minh tương tự, $EH$ (hay $EB$) là phân giác của góc $\widehat{DEF}$.\\
	Xét tam giác $DEF$, $K$ là giao điểm của cạnh đối diện $ED$ kéo dài với $BC$. Theo tính chất chùm điều hoà của tứ giác toàn phần $AB.CD.EF$, bộ điểm $(K, F, B, C) = -1$ là hàng điểm điều hoà.\\
	Vì $O$ là trung điểm của $BC$ nên theo hệ thức Newton cho hàng điểm điều hoà với trung điểm $O$:
	$OB^2 = OC^2 = OF \cdot OK$.\\
	Mặt khác theo hệ thức Maclaurin: $FO \cdot FK = FB \cdot FC$.
	\item Vì $J$ là trung điểm của $AH$ nên $J$ là tâm đường tròn ngoại tiếp tứ giác $ADHE$ (đường kính $AH$).\\
	Trục đẳng phương của đường tròn $(J)$ đường kính $AH$ và đường tròn $(O)$ đường kính $BC$ chính là đường thẳng chứa dây cung chung $ED$ (vì $ED$ kéo dài cắt các đường cao tại các điểm đẳng phương).\\
	Giao điểm $I$ của tiếp tuyến tại $B$ với đường thẳng $ED$ có cùng phương tích đối với các đường tròn liên quan, dẫn đến $I$ nằm trên đường đối cực của $B$.\\
	Do đó đường thẳng $OI$ vuông góc với trục nối tâm hoặc vuông góc với đường thẳng $BJ$, suy ra $OI \perp BJ$.
\end{enumerate}
}
\end{bt}
